\section{Memory, pace and payment}\label{sec:memory}

\subsection{Memory should grow with the pace of change}
A learner in a changing world trades responsiveness against noise: a short memory follows changes but is thrown about by
noise; a long memory averages noise out but lags behind changes. For a single learner the optimal memory grows like
$\tauE^{1/2}$ \citep{lee2024constructing}. The attention network has two candidate memories: the blending weight
$\varepsilon$, with $\taum=1/\ln(1+\varepsilon)$, and self-reinforcement $\kappa$, which also makes the past persist. With a
perfect signal the shortest memory always wins, so E1 adds $\nu=0.5$ and scores runs by log-loss on the true source
(30 agents, $T=0.3$, 15 memories, six paces, four repeats). With $\varepsilon$ as the memory, $\taum^*\propto\tauE^{0.43}$,
approaching $1/2$ from below; with $\kappa$ it barely moves (\cref{fig:e1,tab:e1}). Self-reinforcement is a commitment
knob, not a memory --- which is why it sets the lock-in edge.

\figw{e1}{0.8\linewidth}{Best memory against the pace of change (E1), logarithmic axes.}{e1}

\begin{table}[tbp]\centering\small\caption{E1: optimal memory per pace.}\label{tab:e1}\input{generated/e1}\end{table}

\subsection{When memory sets the pace}
In niche-learning models agents can also change their environment: stabilizers slow the switching they track,
destabilizers speed it up \citep{lee2024constructing}. If a population's memory sets how fast its own world changes,
does the optimal memory settle, cycle or drift? E4 uses the binary environment of \citet{lee2024constructing} with a
mean-field population: for each resident memory, rare invaders of every other memory are tested, and the best response
is iterated (\cref{fig:e4,tab:e4}). In every case iteration converges to a stable memory from both ends, with no cycles or
drift. At that point the memory equals a passive agent's optimum for the environment the population creates, so the
square-root scaling closes on itself. Stabilizers slow their world 15--18-fold. The stable memory is shorter than the one
best for the group (10.7 against about 14 at $\tauE=1000$): shorter-memory invaders free-ride on the calm residents
provide, a quantitative version of the commons argument of \citet{lee2024constructing}.

\figw{e4}{0.8\linewidth}{Best response against the residents' memory, $\tauE=1000$ (E4). Crossings of the dashed identity
are self-consistent memories; all are stable.}{e4}

\begin{table}[tbp]\centering\small\caption{E4: self-consistent memories.}\label{tab:e4}\input{generated/e4}\end{table}

\subsection{Pooling signals}
Agents could pool several rounds of news before each update. Summing distances over the window never helps, because
the source's advantage is the $\rho$ contrast, which summing dilutes. Averaging rewards over a window of $k$ signals does
help: it acts as a sampling window whose optimum grows like $\tauE^{0.42}$, close to the memory scaling (E7,
\cref{fig:e7}).

\figw{e7}{0.8\linewidth}{Log-loss against the pooling window at $\tauE=300$ (E7).}{e7}

\subsection{Paying for attention does not entrench hubs}
If being listened to pays, do hubs use it to stay hubs, as monopoly arguments suggest? In E3 agents earn in proportion to
the attention they relay, and each epoch the worst earner copies the best earner's $\kappa$. Hubs do not entrench: they
change every one or two epochs with or without payment; with a slow source, payment if anything shortens their tenure
(\cref{tab:e3}). The same holds when the next source is drawn in proportion to relayed attention, and on sparse graphs.
Unlike a stabilizer, a hub cannot make itself more reliable by holding attention: nothing in the model lets it change the
source's pace.

\begin{table}[tbp]\centering\small\caption{E3: hub tenure and takeovers over 150 epochs.}\label{tab:e3}\input{generated/e3}\end{table}
